\section{Abstract}\label{sec:1-abstract}
The nano-hertz gravitational wave background recently reported by NANOGrav~\cite{afzal_nanograv_2023} and other pulsar timing arrays (PTAs) collaborations stands as one of the most exciting signals in contemporary physics.
Among the cosmological interpretations, a stochastic gravitational wave background (SGWB) produced by a network of cosmic strings (one-dimensional topological defects formed during symmetry-breaking phase transitions in the early universe~\cite{kibble_topology_1976}) is a well-motivated and observationally relevant candidate.
The microphysics of such a network, encoded in the string tension $G_{\mu}$, the reconnection probability $p$ and the loop distribution model, leaves a characteristic imprint on the spectral energy density $\Omega_{\text{gw}(f)}$.

Traditional Bayesian parameter estimation using ENTERPRICE/PTArcade relies on MCM sampling with an explicit likelihood~\cite{ellis_enterprise_2020,mitridate_ptarcade_2023}, which scales poorly as the model complexity grows and struggles to marginalize efficiently over nuisance parameters such as the competing supermassive black hole binary (SMBHB) background.

In this proposal, we develop a fully open-source, modular \textbf{Simulation-Based Inference (SBI)} pipeline~\cite{cranmer_frontier_2020} targeting the PTA frequency band.
The forward simulator encodes the Velocity-dependent One-Scale (VOS) network evolution equations~\cite{martins_extending_2002} in an expanding FLRW background, computes the resulting SGWB spectral density analytically (via the BOS~\cite{blanco-pillado_number_2014} and LRS~\cite{lorenz_cosmic_2010} loop distribution models), and injects the signal into mock NANOGrav-15ye noise~\cite{johnson_nanograv_2024}.
The inference backend employs Truncated Marginal Neural Ratio Estimation (TMNRE)~\cite{miller_truncated_2021}, implemented via the \texttt{swyft} library~\cite{miller_swyft_2022}, to learn the marginal posteriors over the string microphysics parameters.

This approach is explicitly complementary to the existing \texttt{saqqara} library~\cite{alvey_simulationbased_2024}, which targets LISA and generic signal templates; our contribution is a dedicated PTA-band forward model and a cosmic-string-specific signal module that could ultimately be contributed upstream.

The cluster resources available to the group are utilized for the large-scale parallel simulation campaigns needed to train the neural estimator.


\section{Objectives}\label{sec:1-objectives}
\begin{itemize}
    \item Implement a vectorized, physically complete forward simulator for the SGWB from cosmic string networks in the PTA frequency band, based on the VOS model~\cite{martins_quantitative_1996,martins_extending_2002} and Nambu-Goto loop dynamics;
    \item Develop a modular SBI pipeline (forward simulator + noise model + TMNRE inference) as a reusable open-source software package, designed to be architecturally compatible with the \texttt{saqqara} library~\cite{alvey_simulationbased_2024} for future integration;
    \item Constrain cosmic string microphysics parameters ($G_{\mu}$,~$p$, loop distribution model) using the latest NANOGrav-15yr dataset~\cite{afzal_nanograv_2023};
    \item Benchmark the SBI pipeline against traditional MCMC-based ENTERPRISE/PTArcade results~\cite{ellis_enterprise_2020,mitridate_ptarcade_2023} to validate and quantify the improvement in inference efficiency;
    \item Explore the implicit marginalization capability of TMNRE~\cite{miller_truncated_2021} over the SMBHB contamination as a nuisance source.
\end{itemize}

\section{Methods}\label{sec:1-methods}
\begin{itemize}
    \item The VOS model~\cite{martins_extending_2002} (a pair of coupled ODEs describing the correlation length $L$ and RMS velocity $v$ of the string network in an FLRW background) is solved numerically using a stiff ODE integrator.
    This constitutes the physically relativistic core of the forward mode;
    \item The SGWB spectral energy density $\Omega_{\text{gw}}(f; G_\mu, p)$ is computed for integrating the loop number density (BOS~\cite{blanco-pillado_number_2014}, LRS~\cite{lorenz_cosmic_2010}, or Lorenz-Wands-Martins~\cite{ringeval_cosmological_2007} models) weighted by the GW emission power from cusps and kinks.
    \item The resulting signal is injected into noise model based on the NANOGrav 15-year free-spectrum posterior~\cite{johnson_nanograv_2024}, using the \texttt{enterprise}~\cite{ellis_enterprise_2020} + \texttt{enterprise\_extensions} PTA likelihood frameworks as a reference.
    \item A compression network (1D CNN or UNet architecture, following the \texttt{saqqara} design~\cite{alvey_simulationbased_2024}) is trained to map frequency-domain PTA data to a low-dimensional summary embedding.
    A neural ratio estimator is then trained using \texttt{swyft}~\cite{miller_swyft_2022} to produce marginal posteriors on the parameter vector $\theta = \{\log G_\mu, \log p, \text{loop model index} \}$.
    \item Simulation campaigns are distributed across the group's compute cluster for embarrassingly parallel prior sampling.
    \item Posterior predictive checks and simulation-based calibration (SBC)~\cite{talts_validating_2018} are applied to validate inference quality.
    \item Results are compared to the existing PTArcade/ENTERPRISE analyses of the same dataset~\cite{afzal_nanograv_2023,mitridate_ptarcade_2023} for quantitative benchmarking.

\end{itemize}
